\documentclass[11pt]{article}
\usepackage[margin=1in]{geometry}
\usepackage{amsmath,amssymb,amsthm}
\usepackage{graphicx}
\usepackage[hidelinks]{hyperref}
\setlength{\emergencystretch}{2em}

\newtheorem{theorem}{Theorem}
\newtheorem{lemma}{Lemma}
\newtheorem{corollary}{Corollary}
\newtheorem{remark}{Remark}

\newcommand{\C}{\mathbb C}
\newcommand{\coz}{\operatorname{coz}}
\newcommand{\cl}{\overline}

\title{Pure-State Pi-Base Does Not Force a Separable Right-Support Net}
\author{Run \texttt{fa\_banach\_001}, lane 19}
\date{June 17, 2026}

\begin{document}
\maketitle

\section*{Source Problem and Lane Context}

The source paper is Jinghao Huang, Karimbergen Kudaybergenov, and Rui Liu,
\emph{The ball-covering property of von Neumann algebras and noncommutative
symmetric spaces}, arXiv:2311.14257.  After proving that a commutative
unital \(C^*\)-algebra \(C(K)\) has BCP/UBCP exactly when \(K\) has a
countable \(\pi\)-basis, it asks for a noncommutative \(C^*\)-algebra
characterization:

\begin{center}
\includegraphics[width=0.92\textwidth]{figures/open_problem_crop.png}
\end{center}

The previous lane-19 packet
\[
\texttt{2311.14257\_separable\_right\_support\_net\_ubcp}
\]
proved the following useful sufficient condition: if \(A\) has positive
contractions \(e_n\) such that \(\cl{Ae_n}\) is separable for every \(n\) and
\[
        \|a\|=\sup_n\|ae_n\|
        \qquad(a\in A),
\]
then \(A\) has UBCP.  The present packet shows that pure-state countable
\(\pi\)-base cannot force this support-net condition, even in the
commutative case.

\section*{The Counterexample}

Let
\[
        K=([0,1]\times\{0,1\})\setminus\{(0,0),(1,1)\}
\]
with the lexicographic order:
\[
        (s,i)<(t,j)
        \quad\Longleftrightarrow\quad
        s<t \text{, or } s=t \text{ and } i<j .
\]
Equip \(K\) with the order topology.  This is the standard endpoint-free
split interval, also called the double-arrow compactum.  Let \(A=C(K)\).

\begin{theorem}
\label{thm:main}
The commutative \(C^*\)-algebra \(A=C(K)\) has the following properties:
\begin{enumerate}
    \item its pure state space \(P(A)\cong K\) has a countable \(\pi\)-base;
    \item \(A\) has UBCP;
    \item no nonzero positive \(e\in A\) has separable closed right ideal
          \(\cl{Ae}\).
\end{enumerate}
Consequently, pure-state countable \(\pi\)-base does not imply the existence
of a separable right-support net.
\end{theorem}

\section*{Idea of the Proof}

The endpoint-free split interval is small from the point of view of
\(\pi\)-bases: rational order intervals form a countable \(\pi\)-base.  But
it is locally large from the point of view of separability of function
algebras: every nonempty open subset contains a copy of a Sorgenfrey interval,
hence is not second countable.

For \(C(K)\), the closed principal ideal generated by \(e\) is
\[
        \cl{C(K)e}
        =
        \{f\in C(K):f=0 \text{ on } Z(e)\}
        \cong C_0(\coz(e)).
\]
If this ideal were separable, then \(\coz(e)\) would be second countable.
That is impossible for nonzero \(e\), because \(\coz(e)\) is a nonempty open
subset of the split interval.

\section*{Proof}

\begin{lemma}
\label{lem:pibase}
The split interval \(K\) has a countable \(\pi\)-base.
\end{lemma}

\begin{proof}
For rational \(0\le p<q\le1\), include the order intervals
\[
        ((p,1),(q,0)).
\]
Also include the rational endpoint rays, where the minimum of \(K\) is
\((0,1)\) and the maximum is \((1,0)\),
\[
        [(0,1),(q,0))\quad(0<q\in\mathbb Q)
        \qquad\text{and}\qquad
        ((p,1),(1,0)]\quad(p\in\mathbb Q,\ p<1).
\]
This family is countable.

Let \(W\subset K\) be nonempty open.  By the order-topology basis, \(W\)
contains a nonempty order interval or endpoint ray.  Such a set has positive
width in the first coordinate: otherwise it would be contained in the empty
gap between an adjacent pair \((t,0)<(t,1)\).  Choose rationals \(p<q\) so
that the corresponding rational interval, or one of the endpoint rays, lies
inside that basic open set.  Hence some member of the displayed countable
family is contained in \(W\), proving that the family is a \(\pi\)-base.
\end{proof}

\begin{lemma}
\label{lem:notsecond}
No nonempty open subset of \(K\) is second countable.
\end{lemma}

\begin{proof}
Let \(W\subset K\) be nonempty open.  By Lemma~\ref{lem:pibase}, \(W\)
contains an interval
\[
        ((p,1),(q,0))
\]
with \(p<q\).  Inside this interval, consider the upper copy
\[
        S_{p,q}=\{(t,1):p<t<q\}.
\]
With the subspace topology inherited from \(K\), this is a Sorgenfrey
interval: a neighborhood base at \((t,1)\) is given by sets corresponding to
\[
        [t,s)
        \qquad(t<s<q).
\]

A nontrivial Sorgenfrey interval is not second countable.  Indeed, if
\(\mathcal B\) were a countable base, then for each \(t\in(p,q)\) one could
choose \(B_t\in\mathcal B\) such that
\[
        t\in B_t\subset [t,q).
\]
If \(s<t\), then \(B_t\) cannot equal \(B_s\), because \(B_t\subset[t,q)\)
does not contain \(s\), while \(B_s\) contains \(s\).  Thus
\(\mathcal B\) would contain uncountably many distinct members, a
contradiction.

Therefore \(S_{p,q}\) is not second countable.  Since subspaces of second
countable spaces are second countable, \(W\) itself is not second countable.
\end{proof}

\begin{lemma}
\label{lem:principal}
If \(0\ne e\in C(K)_+\), then \(\cl{C(K)e}\) is not separable.
\end{lemma}

\begin{proof}
Let
\[
        U=\coz(e)=\{x\in K:e(x)\ne0\}.
\]
This is a nonempty open subset of \(K\).  The closed ideal generated by \(e\)
is
\[
        \cl{C(K)e}
        =
        \{f\in C(K): f|_{K\setminus U}=0\},
\]
which is canonically isometric to \(C_0(U)\).

For a locally compact Hausdorff space \(U\), \(C_0(U)\) is separable only if
\(U\) is second countable.  One way to see this is to take a countable dense
set in \(C_0(U)\); the cozero sets of its rational level cuts form a
countable base for \(U\).  But Lemma~\ref{lem:notsecond} says that no
nonempty open subset of \(K\), including \(U\), is second countable.  Hence
\(\cl{C(K)e}\cong C_0(U)\) is not separable.
\end{proof}

\begin{proof}[Proof of Theorem~\ref{thm:main}]
Since \(A=C(K)\) is commutative and unital, its pure state space is
homeomorphic to \(K\), via point evaluations.  Lemma~\ref{lem:pibase} gives a
countable \(\pi\)-base for \(P(A)\).

The source paper's commutative theorem gives UBCP for \(C(K)\), because
\(K\) has a countable \(\pi\)-base.

Finally, Lemma~\ref{lem:principal} says that no nonzero positive \(e\in A\)
has separable \(\cl{Ae}\).  Therefore a separable right-support net cannot
exist: every element of such a net would have to be zero, and then
\(\sup_n\|1\cdot e_n\|=0\ne1\).
\end{proof}

\section*{Consequence for the C*-BCP Program}

This counterexample does not threaten the source theorem.  In fact, it is a
commutative example where the source theorem works perfectly:
\[
        K\text{ has a countable }\pi\text{-base}
        \quad\Longrightarrow\quad
        C(K)\text{ has UBCP}.
\]
What fails is the lane-19 intermediate hope that pure-state countable
\(\pi\)-base should imply a countable separable right-support net.

Thus the right-support-net theorem remains a strong sufficient condition, but
it is not the necessary bridge to the full classification.  A successful
noncommutative pure-state proof must imitate the commutative scalar bump
mechanism more directly.  Equivalently, it needs a weaker local coefficient
or rank-one bump structure, not separability of the whole principal right
ideal \(\cl{Ae}\).

\section*{Novelty and Search Bounds}

Local indexes were searched for arXiv:2311.14257, pure state \(\pi\)-base,
right-support net, split interval, double arrow, separable principal ideal,
and commutative BCP.  The split interval is a standard compact ordered space,
but no previous run packet recorded this obstruction to the right-support-net
bridge.  The packet is therefore a counterexample to an agent-generated
intermediate lemma, not a claim of a new counterexample to the published
source question.

\section*{Human Review Focus}

Reviewers should check:
\begin{enumerate}
    \item the countable \(\pi\)-base for the split interval;
    \item the identification of the upper copy with a Sorgenfrey interval;
    \item the proof that every nonempty open subset of \(K\) is not second
          countable;
    \item the standard commutative ideal identity
          \(\cl{C(K)e}\cong C_0(\coz(e))\);
    \item the distinction between refuting the support-net bridge and
          refuting the source problem.
\end{enumerate}

\begin{thebibliography}{9}

\bibitem{HKL}
Jinghao Huang, Karimbergen Kudaybergenov, and Rui Liu,
\emph{The ball-covering property of von Neumann algebras and noncommutative
symmetric spaces},
arXiv:2311.14257, 2023.

\bibitem{LLLZ}
Minzeng Liu, Rui Liu, Jimeng Lu, and Bentuo Zheng,
\emph{Ball covering property from commutative function spaces to
non-commutative spaces of operators},
arXiv:2111.04921, 2021.

\bibitem{lane19support}
Run \texttt{fa\_banach\_001}, lane 19,
\emph{Separable Right-Support Nets Force UBCP},
solution packet
\texttt{2311.14257\_separable\_right\_support\_net\_ubcp}, 2026.

\end{thebibliography}

\end{document}
